\chapter{Аналіз технологій VPN та методів автентифікації}

Забезпечення захищеного мережевого з'єднання є фундаментальною вимогою сучасних
інформаційних систем. У цьому розділі проводиться системний аналіз технологій
побудови віртуальних приватних мереж, детально розглядається криптографічна
архітектура протоколу WireGuard, досліджуються методи двофакторної автентифікації
та виконується порівняльний аналіз існуючих рішень інтеграції автентифікації
з WireGuard.

%% ============================================================
\section{Огляд технологій віртуальних приватних мереж}
%% ============================================================

Віртуальна приватна мережа (VPN, від англ.\ \textit{Virtual Private Network})~---
це технологія, що забезпечує захищений зашифрований канал передачі даних між
двома або більше вузлами через публічну мережу, зокрема Інтернет. VPN дозволяє
створювати логічно ізольовані мережеві з'єднання, приховуючи реальні IP-адреси
учасників та захищаючи трафік від перехоплення і модифікації~\cite{stallings2019network}.

Основними сценаріями застосування VPN є: захищений віддалений доступ співробітників
до корпоративних ресурсів, об'єднання географічно розподілених офісів
(site-to-site VPN), забезпечення конфіденційності особистих даних при використанні
публічних мереж, а також обхід географічних обмежень у доступі до інтернет-ресурсів.

\subsection{Порівняльний аналіз протоколів VPN}

Найпоширенішими сучасними VPN-протоколами є WireGuard, OpenVPN
та IPsec. Кожен з них має свої архітектурні особливості, переваги та недоліки.

\textbf{WireGuard}~--- сучасний VPN-протокол, розроблений Джейсоном Доненфелдом
та вперше представлений на конференції NDSS 2017~\cite{donenfeld2017wireguard}.
Протокол реалізовано як модуль ядра Linux (починаючи з версії 5.6) і
відзначається надзвичайно малим розміром кодової бази~--- близько 4\,000 рядків
коду (LOC), що суттєво спрощує аудит безпеки. WireGuard використовує фіксований
набір сучасних криптографічних примітивів, що унеможливлює атаки на зниження
рівня криптографічного захисту. Протокол пройшов формальну верифікацію
за допомогою інструменту CryptoVerif~\cite{lipp2019cryptographic, blanchet2019cryptoverif}.

\textbf{OpenVPN}~--- протокол з відкритим вихідним кодом на основі бібліотеки
OpenSSL, доступний з 2002 року. Його кодова база
налічує понад 70\,000 LOC і підтримує широкий спектр криптографічних
алгоритмів, що надає гнучкість, але водночас збільшує поверхню атаки.
Протокол використовує стек TLS для захисту каналу управління та підтримує
роботу як через UDP, так і через TCP~\cite{os3_2020_vpn}.

\textbf{IPsec} (Internet Protocol Security)~--- набір протоколів IETF для
захисту з'єднань на мережевому рівні. IPsec має зрілу багатовендорну підтримку
(Cisco, Juniper, Linux strongSwan), проте характеризується надзвичайно великою складністю
специфікацій~--- згідно з оцінкою~\cite{donenfeld2017wireguard}, кодова база IPsec у сукупності всіх реалізацій сягає понад 300\,000 LOC.
Реалізації IPsec включають IKE (Internet Key Exchange) для обміну ключами
та ESP (Encapsulating Security Payload) або AH (Authentication Header) для
захисту пакетів.

У таблиці~\ref{tab:vpn_comparison} наведено порівняння ключових характеристик
трьох протоколів.

\begin{table}[htbp]
\caption{Порівняння протоколів VPN}
\label{tab:vpn_comparison}
\small
\begin{tabularx}{\textwidth}{|p{4cm}|X|X|X|}
\hline
\textbf{Критерій} & \textbf{WireGuard} & \textbf{OpenVPN} & \textbf{IPsec} \\
\hline
Розмір кодової бази & $\sim$4\,000 LOC & $>$70\,000 LOC & $>$300\,000 LOC \\
\hline
Продуктивність (пропускна здатність) & Висока (у 2--5 разів вища за OpenVPN) & Середня & Висока (апаратне прискорення) \\
\hline
Затримка з'єднання & Мінімальна (1-RTT) & Висока (3--7 RTT) & Середня (IKEv2: 2 RTT) \\
\hline
Набір шифрів & Фіксований (без налаштування) & Конфігурований (ризик хибного вибору) & Конфігурований (ризик хибного вибору) \\
\hline
Формальна верифікація & Так (CryptoVerif, 2019)~\cite{lipp2019cryptographic} & Ні & Часткова \\
\hline
Інтеграція з ядром ОС & Так (Linux 5.6+) & Простір користувача & Так \\
\hline
Складність налаштування & Низька & Середня & Висока \\
\hline
Підтримка UDP/TCP & UDP (основний) & UDP/TCP & UDP \\
\hline
\end{tabularx}
\end{table}

Порівняльний аналіз продуктивності, проведений у роботі~\cite{os3_2020_vpn}
та підтверджений більш пізніми дослідженнями~\cite{mdpi2025wgopenvpn}, свідчить
про те, що WireGuard демонструє у 2--5 разів вищу пропускну здатність порівняно
з OpenVPN в ідентичних умовах тестування. Це досягається завдяки оптимізованій
реалізації в просторі ядра та використанню сучасних криптографічних алгоритмів,
апаратно підтримуваних на більшості сучасних процесорів.

%% ============================================================
\section{Протокол WireGuard: архітектура та криптографічні основи}
%% ============================================================

\subsection{Криптографічні примітиви}

WireGuard будується на чітко визначеному наборі сучасних криптографічних
примітивів, вибір яких обґрунтований з точки зору безпеки та продуктивності~\cite{donenfeld2017wireguard}:

\begin{itemize}
  \item \textbf{Curve25519} (RFC~7748~\cite{rfc7748})~--- алгоритм обміну
        ключами Діффі-Хеллмана на еліптичній кривій, що забезпечує 128-бітний
        рівень безпеки. Використовується для генерації пар ключів (публічний/приватний)
        та у процесі встановлення з'єднання.

  \item \textbf{ChaCha20-Poly1305} (RFC~8439~\cite{rfc8439})~--- аутентифіковане
        шифрування з асоційованими даними (AEAD). ChaCha20 є потоковим шифром,
        а Poly1305~--- MAC-алгоритмом. Поєднання забезпечує конфіденційність,
        цілісність і автентичність даних без необхідності окремих механізмів захисту.

  \item \textbf{BLAKE2s} (RFC~7693~\cite{rfc7693})~--- хеш-функція, оптимізована
        для 32-бітних платформ. Використовується для побудови MAC-кодів та в
        контексті протоколу узгодження ключів.

  \item \textbf{HKDF} (RFC~5869~\cite{rfc5869})~--- функція виведення ключів
        на основі HMAC. Застосовується для отримання симетричних ключів шифрування
        з матеріалу, виробленого у процесі рукостискання.

  \item \textbf{SipHash-2-4}~--- швидка хеш-функція для захисту від HashDoS-атак
        при пошуку вузлів у таблиці маршрутизації.
\end{itemize}

\subsection{Noise Protocol Framework та шаблон рукостискання}

WireGuard використовує протокол рукостискання, побудований на основі
Noise Protocol Framework~\cite{perrin2018noise}~--- загальної специфікації для
побудови захищених протоколів обміну повідомленнями. Конкретний шаблон, що
застосовується WireGuard, позначається як:

\begin{center}
\texttt{Noise\_IKpsk2\_25519\_ChaChaPoly\_BLAKE2s}
\end{center}

Розшифрування цього позначення: \texttt{IK}~--- ініціатор знає статичний
публічний ключ респондента (\textit{Known}), а статичний ключ ініціатора
передається у рукостисканні (\textit{Initiator sends static key});
\texttt{psk2}~--- застосування попередньо узгодженого ключа (PSK) у другому
повідомленні рукостискання (psk2); за відсутності налаштованого PSK використовується
нульовий рядок; \texttt{25519}~--- крива Diffie-Hellman Curve25519;
\texttt{ChaChaPoly}~--- AEAD-шифр ChaCha20-Poly1305; \texttt{BLAKE2s}~--- хеш-функція~\cite{dowling2018provable}.

Рукостискання WireGuard є \textit{1-RTT} (ініціатор надсилає повідомлення,
респондент відповідає~--- загалом два пакети для встановлення сесії). На відміну від TLS~1.2 (2 RTT), WireGuard
досягає меншого часу встановлення з'єднання; порівняно з TLS~1.3, перевага
WireGuard полягає в суттєво простішій машині станів і відсутності переговорів
про набір шифрів. IKEv2 потребує 2 RTT та 4 повідомлення до готовності тунелю.

У процесі рукостискання реалізується:
\begin{enumerate}
  \item взаємна автентифікація сторін за допомогою статичних пар ключів;
  \item узгодження ефемерних ключів для кожної сесії (забезпечує PFS~---
        досконалу пряму секретність);
  \item необов'язкова перевірка PSK для додаткового рівня симетричної
        автентифікації.
\end{enumerate}

На рисунку~\ref{fig:noise_handshake} зображено схему рукостискання WireGuard
за шаблоном Noise IK.

\begin{figure}[htbp]
\centering
\begin{tikzpicture}[
  node distance=5cm,
  box/.style={rectangle, draw, minimum width=2.5cm, minimum height=0.8cm,
              font=\small\bfseries, align=center},
  arr/.style={-{Stealth[length=6pt]}, thick},
  label/.style={font=\small, align=center}
]

%% Вузли-учасники
\node[box] (I) {Ініціатор};
\node[box, right=of I] (R) {Респондент};

%% Вертикальні лінії (lifelines)
\draw[dashed] (I.south) -- ++(0,-7cm) coordinate (Iend);
\draw[dashed] (R.south) -- ++(0,-7cm) coordinate (Rend);

%% Повідомлення 1: Init
\coordinate (m1s) at ($(I.south)+(0,-1.0cm)$);
\coordinate (m1e) at ($(R.south)+(0,-1.0cm)$);
\draw[arr] (m1s) -- (m1e)
  node[midway, above, label]{Handshake Init};
\node[below right, label, xshift=-2cm] at (m1e)
  {$\bigl(E_i,\; \mathrm{DH}(E_i, S_r),\; \mathrm{DH}(S_i, S_r)\bigr)$};

%% Повідомлення 2: Response
\coordinate (m2s) at ($(R.south)+(0,-3.5cm)$);
\coordinate (m2e) at ($(I.south)+(0,-3.5cm)$);
\draw[arr] (m2s) -- (m2e)
  node[midway, above, label]{Handshake Response};
\node[below left, label, xshift=2cm] at (m2e)
  {$\bigl(E_r,\; \mathrm{DH}(E_r, E_i),\; \mathrm{DH}(E_r, S_i)\bigr)$};

%% Тунель
\coordinate (t1s) at ($(I.south)+(0,-5.8cm)$);
\coordinate (t1e) at ($(R.south)+(0,-5.8cm)$);
\draw[arr, double, double distance=2pt] (t1s) -- (t1e)
  node[midway, above, label]{Зашифрований трафік (ChaCha20-Poly1305)};

\end{tikzpicture}
\caption{Схема рукостискання WireGuard (шаблон Noise IK): $E$~--- ефемерний ключ,
$S$~--- статичний ключ, $\mathrm{DH}$~--- операція Diffie-Hellman (Curve25519)
(крок psk2 опущено для наочності)}
\label{fig:noise_handshake}
\end{figure}

Формальна верифікація протоколу за допомогою CryptoVerif підтвердила
виконання властивостей: конфіденційності, автентичності та PFS
за прийнятих моделей загроз~\cite{lipp2019cryptographic, blanchet2019cryptoverif}.

\subsection{Концепція Cryptokey Routing}

Центральним елементом архітектури WireGuard є таблиця маршрутизації за
криптографічними ключами (\textit{Cryptokey Routing Table})~\cite{donenfeld2017wireguard}.
Кожному вузлу WireGuard-інтерфейсу відповідає унікальна пара ключів
(приватний/публічний). Для кожного однорангового вузла (peer) конфігурація
містить:
\begin{itemize}
  \item публічний ключ вузла;
  \item список дозволених IP-адрес (\textit{AllowedIPs})~--- CIDR-підмережі,
        з яких очікуються пакети від цього вузла або на які слід їх направляти;
  \item (необов'язково) зовнішня IP-адреса та порт вузла.
\end{itemize}

При отриманні пакету WireGuard перевіряє, чи відповідає IP-адреса відправника
одному з записів \texttt{AllowedIPs} для будь-якого сконфігурованого вузла.
При відправці~--- маршрутизує пакет до вузла відповідно до таблиці.
Автентифікація пакетів відбувається на рівні криптографічних ключів,
а не IP-адрес чи облікових даних.

%% ============================================================
\section{Проблема автентифікації користувачів у WireGuard}
%% ============================================================

Незважаючи на свої криптографічні переваги, WireGuard принципово не надає
механізмів автентифікації \textit{користувачів}~\cite{donenfeld2017wireguard, wireguard_limitations}.
Протокол ідентифікує вузли виключно за статичними публічними ключами. Це означає:

\begin{itemize}
  \item відсутня перевірка особи (хто саме використовує ключ);
  \item відсутня підтримка двофакторної автентифікації;
  \item неможливе динамічне відкликання доступу без ручного видалення ключа
        з конфігурації на сервері;
  \item відсутній журнал автентифікації (хто і коли підключився).
\end{itemize}

Це є свідомим архітектурним рішенням~\cite{donenfeld2017wireguard, wireguard_limitations}: WireGuard
позиціонується як мінімальний примітив рівня даних (data plane), а управління
ідентичністю та доступом навмисно виведено за межі протоколу. Відтак
побудова повноцінної системи захищеного доступу вимагає додаткового рівня
управління, що і становить предмет цієї дипломної роботи.

%% ============================================================
\section{Методи двофакторної автентифікації}
%% ============================================================

Двофакторна автентифікація (2FA) передбачає перевірку особи користувача
за допомогою двох незалежних факторів, що належать до різних категорій:
знання (пароль), володіння (пристрій, токен) та притаманність (біометрія).
Це суттєво підвищує захист навіть у випадку компрометації одного фактора.

\subsection{HOTP: одноразові паролі на основі HMAC}

Алгоритм HOTP (HMAC-based One-Time Password), стандартизований у RFC~4226~\cite{rfc4226},
генерує одноразові паролі на основі лічильника подій. Формально,
для лічильника $C$ та секретного ключа $K$ значення HOTP обчислюється як:

\begin{equation}
\text{HOTP}(K, C) = \text{Truncate}\bigl(\text{HMAC-SHA-1}(K, C)\bigr)
\label{eq:hotp}
\end{equation}

де функція \textit{Truncate} виділяє 31-бітне значення з 20-байтного виводу
HMAC за фіксованим правилом, яке потім перетворюється на 6-8-значний числовий
пароль за допомогою операції по модулю~$10^d$, де $d$~--- кількість цифр.

Основним недоліком HOTP є необхідність синхронізації лічильника між клієнтом
і сервером: якщо клієнт генерує код, але не використовує його, стан лічильника
розсинхронізується. Стандарт допускає певне вікно розсинхронізації, однак
за значних відхилень потрібна ручна ресинхронізація.

\subsection{TOTP: одноразові паролі на основі часу}

TOTP (Time-based One-Time Password), стандартизований у RFC~6238~\cite{rfc6238},
є розширенням HOTP, де лічильник замінено значенням поточного часу:

\begin{equation}
T = \left\lfloor \frac{T_{\text{current}} - T_0}{X} \right\rfloor
\label{eq:totp_counter}
\end{equation}

\begin{equation}
\text{TOTP}(K, T) = \text{HOTP}(K, T)
\label{eq:totp}
\end{equation}

де $T_{\text{current}}$~--- поточний час в секундах від епохи Unix,
$T_0 = 0$~--- початковий момент часу, $X = 30$~--- розмір часового вікна
(30 секунд). Таким чином, код змінюється кожні 30 секунд.

Перевагою TOTP над HOTP є автоматична синхронізація: клієнту і серверу достатньо
мати синхронізований час (через NTP), немає необхідності зберігати стан лічильника.
Стандарт допускає вікно перевірки $\pm 1$ часовий крок для компенсації мережевих
затримок та невеликих розбіжностей годинників~\cite{rfc6238}.

Алгоритм підтримується застосунками-автентифікаторами (Google Authenticator,
Authy, Microsoft Authenticator) та сумісний зі схемою otpauth:// URI, що дозволяє
передавати секрет через QR-код.

\subsection{WebAuthn/FIDO2: автентифікація на основі асиметричних ключів}

WebAuthn (Web Authentication API), стандартизований W3C у 2019 році~\cite{w3c_webauthn2019},
є частиною специфікації FIDO2 і дозволяє здійснювати автентифікацію за допомогою
апаратних токенів (YubiKey, SoloKey) або вбудованих у пристрій автентифікаторів
(Touch ID, Windows Hello). На відміну від OTP, WebAuthn використовує асиметричну
криптографію: під час реєстрації генерується пара ключів, приватний ключ
зберігається у захищеному сховищі пристрою, публічний~--- на сервері.

WebAuthn є стійким до фішингу, оскільки підпис прив'язаний до origin (домену)
сервера і не може бути повторно використаний на підробленому сайті. Проте
впровадження WebAuthn вимагає браузерного API або спеціальної клієнтської
бібліотеки та підтримки апаратних автентифікаторів.

\subsection{Порівняння методів 2FA}

Таблиця~\ref{tab:2fa_comparison} узагальнює ключові характеристики розглянутих
методів двофакторної автентифікації.

\begin{table}[htbp]
\caption{Порівняння методів двофакторної автентифікації}
\label{tab:2fa_comparison}
\small
\begin{tabularx}{\textwidth}{|p{3.8cm}|X|X|X|}
\hline
\textbf{Критерій} & \textbf{HOTP} & \textbf{TOTP} & \textbf{WebAuthn/FIDO2} \\
\hline
Стандарт & RFC 4226 & RFC 6238 & W3C 2019 \\
\hline
Фактор & Володіння (пристрій) & Володіння (пристрій) & Володіння (токен/пристрій) \\
\hline
Зручність використання & Середня (лічильник) & Висока (автоматична зміна) & Висока (одне підтвердження) \\
\hline
Стійкість до фішингу & Низька & Низька & Висока (прив'язка до origin) \\
\hline
Офлайн-підтримка & Так & Так & Так (апаратний токен) \\
\hline
Потреба у синхронізації & Лічильника & Часу (NTP) & Не потрібна \\
\hline
Складність впровадження & Низька & Низька & Висока (браузерний API) \\
\hline
Залежність від IdP & Ні & Ні & Ні \\
\hline
Підтримка мобільних застосунків & Повна & Повна & Обмежена \\
\hline
\end{tabularx}
\end{table}

Для цілей дипломної роботи TOTP є найбільш практичним вибором: стандарт RFC~6238,
широка підтримка автентифікаторів, відсутність залежності від зовнішніх
провайдерів ідентичності, можливість автономної роботи та низька складність
серверної реалізації.

%% ============================================================
\section{Аналіз існуючих рішень інтеграції автентифікації з WireGuard}
%% ============================================================

Оскільки WireGuard свідомо не включає механізмів автентифікації користувачів,
спільнота розробила кілька підходів до вирішення цієї проблеми.
Нижче розглядаються найбільш поширені відкриті рішення.

\textbf{DefGuard}~\cite{defguard2024}~--- рішення на мові Rust, що реалізує
концепцію «PSK як сесійний токен». Після успішної автентифікації (включаючи 2FA)
сервер динамічно оновлює попередньо узгоджений ключ (PSK) у конфігурації
WireGuard-пристрою. Це розділяє площини управління та даних: без чинного PSK
встановити VPN-з'єднання неможливо. Недоліком є вимога наявності власного клієнта
DefGuard~--- стандартний WireGuard-клієнт не підтримується.

\textbf{Firezone}~\cite{firezone2024}~--- система з управляючою площиною на Elixir
та площиною даних на Rust. Використовує OIDC (OpenID Connect) для автентифікації:
після перевірки особи через OIDC-провайдер (Google, Okta тощо) сервер видає
публічний ключ WireGuard клієнту. MFA делегується OIDC-провайдеру.
Потребує зовнішнього постачальника ідентичності.

\textbf{NetBird}~\cite{netbird2024}~--- рішення для побудови повносітчастих
оверлейних мереж (full mesh overlay) на мові Go. Автентифікація та MFA
делегуються зовнішньому провайдеру ідентичності через OIDC.
NetBird більше орієнтований на сценарій «zero-trust network» для корпоративних
середовищ, ніж на класичну hub-and-spoke VPN-архітектуру.

\textbf{Tailscale / Headscale}~\cite{tailscale2024,headscale2024}~--- модель
координаційного сервера для управління ключами вузлів. Tailscale є комерційним
рішенням, Headscale~--- його відкрита самостійна реалізація. MFA делегується
зовнішньому IdP (Google, Microsoft, GitHub). Потрібен клієнт Tailscale
(відкритий, але непов'язаний з \texttt{wg-quick}); стандартний
\texttt{wg-quick}-клієнт не підтримується.

\textbf{WAG (WireGuard Authentication Gateway)}~\cite{wag2024}~--- легке
рішення, що реалізує примусову TOTP-автентифікацію на маршрутному рівні.
До успішної 2FA-верифікації маршрути до захищених ресурсів є недоступними.
Підтримує стандартний WireGuard-клієнт. Однак рішення не надає управління
ключами та орієнтоване на примусову перевірку на шлюзі.

\subsection{Порівняння існуючих рішень}

У таблиці~\ref{tab:solutions_comparison} наведено зведений порівняльний аналіз
розглянутих рішень.

\begin{table}[htbp]
\caption{Порівняння існуючих рішень інтеграції автентифікації з WireGuard}
\label{tab:solutions_comparison}
\small
\begin{tabularx}{\textwidth}{|p{2.5cm}|X|X|X|X|X|}
\hline
\textbf{Критерій} & \textbf{DefGuard} & \textbf{Firezone} & \textbf{NetBird} & \textbf{Tailscale/ Headscale} & \textbf{WAG} \\
\hline
Механізм 2FA & PSK-токен + TOTP/FIDO & OIDC (MFA у IdP) & OIDC (MFA у IdP) & OIDC (MFA у IdP) & TOTP (вбудований) \\
\hline
Точка примусу & Управляюча площина & Видача ключів & Видача ключів & Видача ключів & Маршрутизація \\
\hline
Стандартний WG-клієнт & Ні (власний) & Так & Ні (власний) & Ні (власний) & Так \\
\hline
Self-hosted & Так & Так & Так & Headscale: Так & Так \\
\hline
Залежність від IdP & Ні & Так & Так & Так & Ні \\
\hline
Стек технологій & Rust & Elixir/Rust & Go & Go & Go \\
\hline
Управління ключами & Так & Так & Так & Так & Ні \\
\hline
\end{tabularx}
\end{table}

Ключовим спостереженням є те, що жодне з розглянутих рішень не задовольняє
водночас усі наступні вимоги:
\begin{itemize}
  \item підтримку стандартного WireGuard-клієнта;
  \item відсутність залежності від зовнішнього IdP;
  \item власний механізм TOTP;
  \item повноцінне управління ключами.
\end{itemize}
Крім того, у науковій літературі
відсутні рецензовані публікації, присвячені систематичному підходу до
інтеграції 2FA з WireGuard~--- наявні лише практичні реалізації у відкритому
програмному забезпеченні.

%% ============================================================
\section{Висновки до розділу~1}
\label{sec:chapter1_conclusions}
%% ============================================================

За результатами проведеного аналізу можна сформулювати такі висновки:

По-перше, протокол WireGuard є оптимальним вибором для реалізації тунельного
рівня VPN-системи. Його переваги~--- мала кодова база (4\,000 LOC), формально
верифікована криптографія~\cite{lipp2019cryptographic}, фіксований захищений
набір шифрів (відсутність атак на зниження рівня), інтеграція з ядром Linux та
вища продуктивність порівняно з OpenVPN~\cite{mdpi2025wgopenvpn}~---
роблять його технічно обґрунтованим підґрунтям для системи.

По-друге, алгоритм TOTP (RFC~6238)~\cite{rfc6238} є найбільш практичним
методом двофакторної автентифікації для розроблюваної системи. TOTP не потребує
зовнішнього постачальника ідентичності, є самодостатнім стандартом, широко
підтримується мобільними автентифікаторами та має низьку складність
серверної реалізації. На відміну від WebAuthn, він не вимагає спеціального
апаратного забезпечення.

По-третє, аналіз існуючих рішень виявив суттєву прогалину: жодне з них не
поєднує підтримку стандартного WireGuard-клієнта з власним TOTP-механізмом
та повноцінним управлінням ключами без залежності від зовнішнього IdP.
Академічна спільнота також не запропонувала систематичного підходу до вирішення
цієї задачі.

Отже, розробка власної системи VPN з двофакторною автентифікацією
є доцільною як з практичної, так і з наукової точки зору. Розроблювана система
повинна: (1)~використовувати WireGuard як транспортний рівень; (2)~реалізувати
TOTP-автентифікацію без залежності від зовнішніх сервісів; (3)~забезпечити
управління ключами WireGuard через захищений API; (4)~підтримувати стандартний
WireGuard-клієнт для зручності розгортання; (5)~надавати механізм динамічного
відкликання доступу.


\subsection*{Подальші напрямки досліджень}

Окремим напрямком подальшого розвитку є постквантовий WireGuard~\cite{hulsing2021postquantum},
що пропонує заміну алгоритмів Curve25519 на постквантові аналоги для захисту
від загроз з боку квантових обчислень, проте виходить за межі цієї роботи.

\newpage
